\subsection{An actual full gentle witness in the kernel of the CM character}

The following argument supplies the missing strictness clause of
historical A.7.  Its local starting polynomial is integral, and every
subsequent lift is an automorphism of an actual compact group.
The formal polynomial is not being substituted for that lifting step.

\begin{result}{E.197}{Proved}
{An integral degree-seven invariant has a full gentle realization}
Let \(\mathsf L(w)\) be the standard Lyndon bracket of a word in
\(x<y\), with the longest proper Lyndon suffix used at each step.
Put
\[
\begin{split}
h={}&-3\mathsf L(xxxyxyy)-3\mathsf L(xxxyyxy)
       -\mathsf L(xxyxyxy)\\
   &+3\mathsf L(xxyxyyy)+3\mathsf L(xxyyxyy)
       +3\mathsf L(xxyyyxy)+\mathsf L(xyxyxyy).
\end{split}
\tag{E.197.1}
\]
There is a full profinite gentle pair
\(g=(1,f)\in\widehat{\mathrm{GT}}_0\) whose image in the
free pro-\(11\) group \(P\) satisfies
\[
 f_{11}\in\gamma_7P,\qquad
 f_{11}\bmod\gamma_8P=h
       \quad\hbox{in }L_7(\mathbb Z_{11}x\oplus\mathbb Z_{11}y).
\tag{E.197.2}
\]
Its image in every other maximal pro-prime quotient is the identity.
Here \(\gamma_i\) denotes the closed ordinary lower central series.
No order or finite-level-only condition replaces full membership in
\(\widehat{\mathrm{GT}}_0\).
\end{result}

\begin{proof}
\emph{The integral identities and their actual outer meaning.}
In this paragraph \(x,y,z=-x-y\) are Lie elements.  Define
\(k=-h(z,x)\).  Expansion in the free associative algebra over
\(\mathbb Z\) gives the four identities
\[
\begin{aligned}
h(x,y)+h(y,x)&=0,\\
[x,h(x,y)]&=[x+y,h(z,y)],\\
[h(x,y),y]&=[k,x+y],\\
h(x,y)+h(y,z)+h(z,x)&=0.
\end{aligned}
\tag{E.197.3}
\]
The reproducible coefficient certificate below checks the entire
identities, not selected evaluations.  The polynomial \(h\) has
62 nonzero associative coefficients.  Write
\(D_h(x)=0\), \(D_h(y)=[h,y]\).  The third identity says
\(D_h(z)=[k,z]\), so this is a special derivation.

Use the literal modular marking
\[
\mathcal B=\widehat{C_2*C_3}
 =\langle s,r\mid s^2=r^3=1\rangle^{\wedge},\quad
\tau=sr,\quad F=\ker(\mathcal B\to\mathfrak S_3),
\quad x=\tau^{-2},\quad y=sxs.
\tag{E.197.4}
\]
The six-coset Schreier calculation gives the actual free profinite
basis \(x,y\) of \(F\).  In the group put \(z=(xy)^{-1}\).
Reduction of words in the free product gives
\[
 r:(x,y,z)\longmapsto(y,z,x),\qquad
 s:(x,y,z)\longmapsto(y,x,x^{-1}zx).
\tag{E.197.5}
\]
These are the same marking and conjugations used in E.180.
On the associated Lie algebra their linear actions are
\(s(x,y)=(y,x)\), \(r(x,y)=(y,-x-y)\).
The first two identities in E.197.3 now imply, as equalities of
derivations,
\[
sD_hs^{-1}-D_h=-\operatorname{ad}h,
\qquad rD_hr^{-1}-D_h=-\operatorname{ad}h.
\tag{E.197.6}
\]
For the second identity the values at \(x,y\) on the left are
\(-[h(y,z),z]\) and \(-[h,y]\); the second and first identities
in E.197.3 identify the former with \(-[h,x]\).
Thus both \([D_h]\) and \(-[D_h]\) are invariant outer
derivations for the actual geometric group of order six.

\emph{Integral realization before taking fixed points.}
For a free pro-\(11\) group the integral Magnus and Hall
coordinates identify
\(\gamma_jP/\gamma_{j+1}P\) with the ordinary free Lie
lattice \(L_j(\mathbb Z_{11}^2)\).  One way to see that no
extra completion relation is introduced is to collect in every
finite nilpotent quotient of the free group: the unique Hall
coordinates complete from \(\mathbb Z\) to \(\mathbb Z_{11}\),
and the leading commutator formula is the integral Lie bracket.
The successive free nilpotent quotients separate the original
free pro-\(11\) group, since every finite \(11\)-group is
nilpotent.  We use these ordinary lower central coordinates,
not the restricted Lie algebra of a different filtration.

Here is the needed onto realization argument.  Suppose a degree
\(n\ge2\) derivation fixes \(x\), has values \([u,y]\)
at \(y\) and \([v,z]\) at \(z\), and preserves
\(x+y+z=0\).  Choose group elements \(U,V\in\gamma_nP\)
with those leading terms.  The prospective images are
\[
             x,\qquad UyU^{-1},\qquad VzV^{-1}.
\]
Their product error has zero term in degree \(n+1\), so lies
in \(\gamma_{n+2}P\).  At any subsequent degree \(m\),
altering \(U,V\) in degree \(m-1\) changes that error by
\[
(a,b)\longmapsto[a,y]+[b,-x-y].
\tag{E.197.7}
\]
This map is onto \(L_m(\mathbb Z_{11}^2)\).  Indeed its image
is \([L_{m-1},x]+[L_{m-1},y]\), which is all of \(L_m\):
the Jacobi identity moves a generator to the outer entry of
every bracket monomial.  Cancel the error successively in every
degree.  The changes converge in \(P\), because a finite
\(11\)-group quotient kills every sufficiently deep
\(\gamma_jP\).  Their limiting product is exactly one.
The images of \(x,y\) define an endomorphism preserving all
three peripheral conjugacy classes.  Its Frattini action is the
identity, so it is onto and hence an automorphism of the finitely
generated pro-\(11\) group.  Applying this with \(u=-h\),
\(v=-k\) realizes the leading derivation \(-D_h\).

Let \(\mathcal S\) be the closed group of special outer
automorphisms of \(P\): its representatives preserve the three
individual peripheral conjugacy classes with exponent one.
It is pro-\(11\).  In fact its representatives have trivial
Frattini action; the successive kernels of their actions on
finite characteristic lower exponent-\(11\) quotients are
\(11\)-groups, and passing to the quotient by inner
automorphisms preserves this assertion.  Define
\[
 I^n=\ker\bigl(\mathcal S\longrightarrow
                   \operatorname{Out}(P/\gamma_{n+1}P)\bigr).
\]
These are closed pro-\(11\) groups, invariant under the actual
\(\mathfrak S_3\) action.  Their leading map identifies
\(I^n/I^{n+1}\) with its image in degree-\(n\) derivations
modulo inner derivations.  To check this description, adjust a
representative by a lift of its inner action modulo
\(\gamma_{n+1}\); its errors on generators then lie in
\(\gamma_{n+1}\).  Changing that adjustment changes the leading
derivation by an inner degree-\(n\) derivation.  A zero outer
leading class is exactly the next kernel \(I^{n+1}\).
The preceding explicit construction shows that
\(-[D_h]\) belongs to this actual image at \(n=7\).

For completeness, fixed points lift across an onto equivariant
homomorphism of pro-\(11\) groups when the acting group has
order prime to eleven.  If \(a\) lifts a fixed target element,
its discrepancy \(a^{-1}\sigma(a)\) is a cocycle in the
kernel.  The two corresponding complements in the semidirect
product are conjugate by that kernel.  For a finite
\(11\)-kernel this follows by induction through an invariant
central elementary abelian subgroup: averaging, with the order
of the acting group invertible, kills \(H^1\) of that subgroup;
induction deals with the quotient.  For a compact pro-\(11\)
kernel take finite invariant quotients and use compactness of
the closed sets of conjugating elements.  Apply this assertion
to \(I^7\to I^7/I^8\) and E.197.6.  It gives an actual
\(\beta\in I^7\) fixed by the full \(\mathfrak S_3\),
with leading class precisely \(-[D_h]\).

\emph{Recovering the literal modular automorphism and the pair.}
Put \(B_{11}=\mathcal B/O^{11}(F)\), so its kernel over
\(\mathfrak S_3\) is \(P\).  The free pro-prime centralizer
calculation gives \(C_P(x)=\langle x\rangle\) and
\(C_P(y)=\langle y\rangle\), and in particular \(Z(P)=1\).
For example, after removing the abelian power of \(x\), a
first nonzero Magnus term of an element centralizing \(x\)
would be a Lie element of degree at least two commuting with
\(x\).  In the free associative algebra such a homogeneous
tensor is a scalar multiple of a pure power of \(x\), which
has zero intersection with that Lie degree.  This proves the
claim degree by degree.  The same argument is given integrally
in \citet[Lemma~2.2, p.~368]{SymCohenPrassidis2006}.

Write \(\omega:\mathfrak S_3\to\operatorname{Out}(P)\)
for the outer action.  There is an actual isomorphism
\[
B_{11}\ \xrightarrow{\ \sim\ }\ 
\operatorname{Aut}(P)
             \mathbin{\times_{\operatorname{Out}(P)}}
                          \mathfrak S_3,
\qquad b\longmapsto(\operatorname{Ad}_b,\bar b).
\tag{E.197.8}
\]
Its kernel is \(Z(P)\).  Given a pair in the target, choose
\(b\) over its second coordinate; the first coordinate differs
from \(\operatorname{Ad}_b\) by an inner automorphism of
\(P\), yielding a preimage.  Compactness makes the resulting
bijection a topological isomorphism.
Choose a representative \(\psi\) of \(\beta\) fixing
\(x\) literally, using its peripheral conjugator at \(x\).
Since \(\beta\) commutes with \(\omega(\mathfrak S_3)\),
\((a,d)\mapsto(\psi a\psi^{-1},d)\) preserves E.197.8.
It gives an actual automorphism \(\phi\) of \(B_{11}\),
restricting to \(\psi\).

We have \(\phi(\tau)^2=\phi(x^{-1})=\tau^2\).
Both \(\phi(\tau)\) and \(\tau\) centralize \(x\), and
their quotient belongs to \(P\); thus
\(\phi(\tau)=x^a\tau\).  Squaring gives \(2a=0\) in
\(\mathbb Z_{11}\), so \(\phi(\tau)=\tau\).
Let \(f_1=s\phi(s)\in P\).  From \(\phi(s)^2=1\)
and the swap action of \(s\), its abelianization is
\((a,-a)\).  Replace \(\phi\) by
\(\phi_0=\operatorname{Ad}(x^a)\phi\).  This still fixes
\(\tau\), and
\[
 f_{11}=s\phi_0(s)=y^a f_1x^{-a}\in P',\qquad
 \phi_0(s)=sf_{11},\qquad \phi_0(r)=f_{11}^{-1}r.
\tag{E.197.9}
\]
Consequently its restriction to \(P\) is exactly
\(x\mapsto x,\ y\mapsto f_{11}^{-1}yf_{11}\).
The relations \(s^2=r^3=1\), using E.197.5, give
\[
 f_{11}(y,x)f_{11}(x,y)=1,\qquad
 f_{11}(z,x)f_{11}(y,z)f_{11}(x,y)=1.
\tag{E.197.10}
\]
These are identities in the full pro-\(11\) group.

There is no uncontrolled change to the leading coefficient in
this normalization.  For \(j\ge2\), the same Magnus argument
proves
\[
C_{P/\gamma_{j+1}P}(x)
          =\langle x\rangle\gamma_jP/\gamma_{j+1}P,
\tag{E.197.11}
\]
and similarly at \(y\).  Since \(\beta\in I^7\), the
normalized restriction is \(\operatorname{Ad}(w)\) modulo
\(\gamma_8P\).  It fixes \(x\), so
\(w\in\langle x\rangle\gamma_7P\) there.  Its value at
\(y\) gives \(f_{11}w\in\langle y\rangle\gamma_7P\).
The abelianization of \(f_{11}\) is zero, forcing both powers
to vanish.  Hence \(f_{11}\in\gamma_7P\).  If its leading
term is \(h_f\), its normalized leading derivation is
\(-D_{h_f}\).  Its difference from \(-D_h\) is inner and
vanishes at \(x\).  A degree-seven inner element commuting
with \(x\) is zero; evaluation at \(y\) then gives
\(h_f=h\).  This proves E.197.2 before the full lift.

\emph{The full profinite lift with its actual source retained.}
The canonical comparison is
\[
\mathcal B/R_{\rm nil}(F)
       \ \simeq\ 
       \prod_{p,\mathfrak S_3}\mathcal B/O^p(F).
\tag{E.197.12}
\]
Indeed \(F/R_{\rm nil}(F)\) is its maximal pronilpotent
quotient and hence the product of its maximal pro-prime
quotients.  At every finite level the distinct prime-power
kernels have no common nontrivial quotient, so the map onto
the displayed fibre product is onto; compactness gives the
full statement.  The outer quotient is the same
\(\mathfrak S_3\) in every factor.
Put \(\phi_0\) in the eleven-primary factor and the identity
elsewhere.  All induce the identity on \(\mathfrak S_3\).
The new and old positive generators \(s,r\) in \(B_{11}\)
are conjugate by elements of \(P\): they give complements
over the cyclic groups of orders two and three, so the coprime
complement argument above applies.  Those conjugators, with
identity entries at the other primes, are actual elements of
E.197.12.  The conjugator at \(\tau\) is literally one.
Thus this is a multiplier-one peripheral automorphism with
three specified conjugators on the actual pronilpotent source.

Apply the independently proved central lifting criterion E.57
with \(N=R_{\rm nil}(F)\).  Its final commutator equality
\(\overline{[N,\mathcal B]}=N\) supplies a full peripheral
automorphism of \(\mathcal B\), retaining all three
conjugators modulo \(N\).  Its conjugator at \(\tau\) lies
in \(N\), because that prescribed quotient conjugator is one.
Conjugate this full automorphism by the inverse of that element.
The resulting \(\Phi\) fixes \(\tau\) exactly and has the
same action and all the same conjugators modulo \(N\).
It induces
the identity on \(\mathfrak S_3\) and preserves \(F\).
Set \(f=s\Phi(s)\).  Its eleven-primary image is the
specified \(f_{11}\); all other pro-prime images are one.
In particular its image in
\(F^{\rm ab}=\prod_p\mathbb Z_p^2\) vanishes, so \(f\in F'\).
The identical word calculation gives
\[
\begin{aligned}
\Phi(s)&=sf,& \Phi(r)&=f^{-1}r,\\
\Phi(x)&=x,& \Phi(y)&=f^{-1}yf,\\
f(y,x)f(x,y)&=1,& f(z,x)f(y,z)f(x,y)&=1.
\end{aligned}
\tag{E.197.13}
\]
The pure restriction is an actual automorphism.  These are
exactly the multiplier-one gentle conditions in the fixed
free profinite marking.  Hence \((1,f)\) is a full gentle
element with all claimed coordinates.  This use of E.57 is
ordinary lifting only: no involution, equivariance of a full
lift, or additional norm equation has been requested.
\end{proof}

\begin{result}{E.198}{Proved}
{The historical CM kernel strictly contains the full GT group}
For the original signed character of E.95 one has
\[
\widehat{\mathrm{GT}}\ \subsetneq\ker\Delta
             \ \subsetneq\widehat{\mathrm{GT}}_0.
\tag{E.198.1}
\]
In particular the entire strict-inclusion and vanishing assertion
of historical GT A.7 holds.  A witness for the first strict
inclusion is the square of the actual full element of E.197.
Its pentagon defect is detected in an actual finite
eleven-primary nilpotent quotient of the five-point group.
\end{result}

\begin{proof}
Write \(t_{ij}\) for the degree-one infinitesimal pure braid
generators.  The degree-seven pentagon of E.197.1 is
\[
\begin{split}
\mathfrak p(h)={}&h(t_{12},t_{23}+t_{24})
                 +h(t_{13}+t_{23},t_{34})-h(t_{23},t_{34})\\
 &-h(t_{12}+t_{13},t_{24}+t_{34})-h(t_{12},t_{23}).
\end{split}
\tag{E.198.2}
\]
There is a Lie homomorphism from the integral infinitesimal
pure braid algebra to the commutator Lie algebra of
\(\mathbb Z[\mathfrak S_4]\), given by
\(t_{ij}\mapsto(ij)\).  Disjoint transpositions commute;
for three distinct indices direct multiplication gives
\([(ij),(ik)+(jk)]=0\).  These are exactly the defining
relations.  The complete coefficient calculation gives
\[
              \mathfrak p(h)\longmapsto
                          160\bigl((34)-(12)\bigr).
\tag{E.198.3}
\]
It is nonzero even modulo eleven.  The code checks all the
infinitesimal braid relations as well as this entire group-ring
identity, so the evaluation is a valid nonvanishing certificate.

One can check the literal cyclic sphere convention directly as well.
The five-cycle in \citet[Definition~1.1, p.~657]{LitSchneps2006Automorphisms}
has leading sum
\[
h(t_{34},t_{45})+h(t_{51},t_{12})+h(t_{23},t_{34})
             +h(t_{45},t_{51})+h(t_{12},t_{23}).
\tag{E.198.cyclic}
\]
Set \(C=\sum_{1\le i<j\le4}(ij)\), a central group-ring
element.  Over \(\mathbb Z_{11}\), the assignment
\(t_{ij}\mapsto(ij)-C/6\) annihilates the full-twist sum.
The sphere relations give
\(t_{i5}=-\sum_{j\le4,j\ne i}t_{ij}\).
Central shifts disappear from every evaluated Lie monomial of
degree at least two.  Substitution into E.198.cyclic gives precisely
\(-160((34)-(12))\), also checked by the certificate.
Thus the exact cyclic convention has the same nonvanishing
detector, with the displayed harmless overall sign.

We explain the passage to an actual finite target.  The integral
graded Lie algebra of the ordinary pure braid group \(PB_4\)
is the Lie algebra on \(t_{ij}\) with exactly the above
relations.  Its homogeneous groups are free abelian.  Both the
presentation and the split free-Lie decomposition are given in
\citet[\S2, pp.~366--367]{SymCohenPrassidis2006}; the latter
decomposes the additive graded groups as those of free groups
of ranks one, two and three.  The full twist is central and
has primitive abelian vector with every coordinate one.
The linking character at the first pair therefore splits it
off as a direct central \(\mathbb Z\) factor.  Passing to
\(\Gamma_{0,5}=PB_4/Z(PB_4)\) removes only this central
degree-one summand; it does not remove E.198.3 in degree seven.

For clarity, the required completion statement can be obtained
directly at this fixed class, without identifying an unspecified
formal deformation with a profinite automorphism.  Put
\(Q=\Gamma_{0,5}/\gamma_8\Gamma_{0,5}\).
Its central factors are free abelian, so an ordered basis of
them gives unique integral collecting coordinates.  Hall
collection gives multiplication and inverse as rational
polynomials in these coordinates, with denominators whose
prime factors are at most seven.  This follows either by the
usual repeated commutator collection, whose coefficients are
products of \(\binom{a}{j}\), \(j\le7\), or from the
fixed-class Hall calculation in E.32.  All denominators are
units in \(\mathbb Z_{11}\).  Completing the coordinates
therefore defines a compact nilpotent group on
\(\mathbb Z_{11}^{d}\); the group identities hold by density
of the integral coordinates.  Reduction of every coordinate
modulo \(11^a\) is an actual finite \(11\)-group quotient,
and their inverse limit retains each lower central coordinate.
Thus the degree-seven group in this actual pro-\(11\)
completion is the integral degree-seven Lie lattice tensored
with \(\mathbb Z_{11}\).  Here the Witt dimensions give
\(d=5+4+10+21+54+125+330=549\); no dimension is needed
for the nonvanishing implication.

For an element of \(\gamma_7P\), substitution into two
elements of this class-seven group depends only on their
degree-one images and on its leading Lie polynomial.
Higher terms vanish, and the leading terms of products add.
The literal five insertion maps defining the group pentagon
therefore have degree-seven defect E.198.2.  By E.197, the
full pair \(g\) has exactly the required eleven-primary
image.  Equation E.198.3 proves that its full pentagon is
nontrivial in an actual finite quotient (already after
coordinate reduction modulo eleven).

Let \(g_0=g^2\), with the gentle group law.  Since
\(\Delta\) is a homomorphism to \(\mathbb Z/2\mathbb Z\),
\(\Delta(g_0)=0\).  The eleven-primary leading term of its
pair is \(2h\): the normalized automorphism for \(g\)
is identity on \(\gamma_7P/\gamma_8P\), and composition
therefore adds the two leading conjugators.  The pentagon
defect of \(g_0\) has group-ring image
\[
                 320\bigl((34)-(12)\bigr),
\tag{E.198.4}
\]
whose two nonzero coefficients are units modulo eleven.
Consequently \(g_0\notin\widehat{\mathrm{GT}}\), while
\(g_0\in\ker\Delta\).  The inclusion on the left of
E.198.1 is E.191.  The strict inclusion on the right is
E.97, which supplies an actual full gentle element of
character one.  All three groups and both witnesses in
E.198.1 are full profinite groups and elements.  Together
with the standard faithful Galois inclusion recalled in
\citet[\S1, p.~657]{LitSchneps2006Automorphisms}, this gives
the entire displayed chain in historical A.7; no assertion
about surjectivity of that Galois inclusion is made.
\end{proof}

\paragraph{Reproducible certificate and audit scope.}
The standalone program \path{code/check_gt_strict_inclusion.py}
expands the seven Lyndon brackets over \(\mathbb Z\), checks
E.197.3 and both generator values in E.197.6, verifies the
literal free-product conjugations E.197.5, checks the
infinitesimal braid relations in \(\mathbb Z[\mathfrak S_4]\),
and evaluates the whole pentagon.  Its JSON companion retains
the complete coefficient data and the SHA256 digest
\[
\begin{gathered}
\texttt{de390b3b14e063e6ea168c52f52f9a0af7}\\[-2pt]
\texttt{c6066d4cc148855444a490a9c11dba}.
\end{gathered}
\]
This is an independently reconstructed replacement for the
unavailable historical Part B certificate; no theorem from
that unrefereed text is used as an input.  The use of the
recovered LM7 proof as a search lead is recorded in the notes.
E.198 proves the exact qualitative strictness assertion;
it does not identify every original finite relative witness
or supply the original numerical compiler of A.4.57.
